\section{Third-law template and ``heat death'': the golden attractor}
\label{sec:third_law}

ASM makes the ``heat death'' problem concrete: if scan slopes are rational (or effectively rational at the accessible resolution), the orbit locks into a short cycle, collapsing the accessible macrostate space. Conversely, badly approximable slopes resist such locking.

\subsection{Rational locking as periodic thermal death}

For $\alpha=p/q$, the orbit visits only $q$ phase points. Under any fixed finite-resolution window family, the induced readout becomes eventually periodic with period dividing $q$. The accumulated mismatch grows linearly (Proposition~\ref{prop:rational_linear}), so a fixed fraction of readout steps become irreparably ``miscounted'' with respect to the uniform reference. Operationally, this is a form of crystalline equilibrium: the system loses macroscopic openness by phase locking.

\subsection{Why the golden branch is special}

The mechanism preventing rapid periodic lock-in is arithmetic. Irrational slopes are approximated by rationals via continued fractions; small denominators yield short pseudo-periods. A slope is \emph{badly approximable} if it cannot be too well approximated by rationals, equivalently if its continued-fraction coefficients are bounded \cite{Cassels1957}. Such slopes yield strong uniform-distribution control and hence logarithmic mismatch bounds (Theorem~\ref{thm:log_bound}).

The golden branch $\alpha=\varphi^{-1}=[0;1,1,1,\dots]$ is extremal in the classical Hurwitz/Markov sense. Hurwitz's theorem states that for any irrational $\alpha$ there exist infinitely many rationals $p/q$ such that
\begin{equation}
  \left|\alpha-\frac{p}{q}\right| < \frac{1}{\sqrt{5}\,q^2},
\end{equation}
and the constant $\sqrt{5}$ is optimal: for $\alpha=\varphi^{-1}$ one has the \emph{uniform} lower bound
\begin{equation}
  \left|\varphi^{-1}-\frac{p}{q}\right| \ge \frac{1}{\sqrt{5}\,q^2}\qquad\text{for all } \frac{p}{q}\in\QQ,
\end{equation}
so $\varphi^{-1}$ is among the hardest irrationals to approximate by rationals \cite{Cassels1957}. Operationally, this delays the onset of short pseudo-periods at a given approximation tolerance, maximizing resistance to phase-locking.

In ASM this means:
\begin{itemize}
  \item long avoidance of short pseudo-periods (maximal resistance to phase locking),
  \item simultaneously, nonzero but controlled mismatch (entropy production remains finite and computable).
\end{itemize}
This motivates a third-law template: perfect ``zero-friction'' readout (vanishing mismatch at all depths) is unattainable for finite observers; attempting to suppress mismatch indefinitely drives the system toward periodic lock-in (loss of openness), rather than toward a frictionless continuum.

\subsection{A structural route to \texorpdfstring{$1/f$}{1/f}-type noise}

Many complex systems exhibit $1/f$-type spectra, often associated with hierarchical time scales and near-criticality \cite{DuttaHorn1981,Weissman1988}. In ASM the Zeckendorf/Fibonacci hierarchy provides a canonical ladder of time scales. When a readout protocol aggregates approximately equal-weight relaxation contributions across this ladder, $1/f$ behavior follows quantitatively (Appendix~\ref{app:fibonacci_1f}).
\noindent In the minimal closure where readout aggregates approximately equal-weight relaxation contributions across a Fibonacci (asymptotically geometric) ladder of characteristic times, one obtains a quantitative $1/f$ band with an explicit prefactor (Appendix~\ref{app:fibonacci_1f}, Proposition~\ref{prop:1overf}):
\begin{equation}
  S(f)\ \approx\ \frac{w}{4 f\,\log\varphi}
\end{equation}
over the intermediate frequency range separating the smallest and largest Zeckendorf-resolved scales.
More generally, bounded nonuniform ladder weights still yield a robust $1/f$ mid-band; the weights affect only the prefactor, which is bracketed between $w_-/(4 f\log\varphi)$ and $w_+/(4 f\log\varphi)$ (Appendix~\ref{app:fibonacci_1f}, Corollary~\ref{cor:1overf_robust}).


